% concatenated from document.tex
%% Copyright 2007-2020 Elsevier Ltd
%% 
%% This file is part of the 'Elsarticle Bundle'.
%% ---------------------------------------------
%% 
%% It may be distributed under the conditions of the LaTeX Project Public
%% License, either version 1.2 of this license or (at your option) any
%% later version.  The latest version of this license is in
%%    http://www.latex-project.org/lppl.txt
%% and version 1.2 or later is part of all distributions of LaTeX
%% version 1999/12/01 or later.
%% 
%% The list of all files belonging to the 'Elsarticle Bundle' is
%% given in the file `manifest.txt'.
%% 
%% Template article for Elsevier's document class `elsarticle'
%% with harvard style bibliographic references

\documentclass[preprint,12pt]{elsarticle}

%% Use the option review to obtain double line spacing
%% \documentclass[authoryear,preprint,review,12pt]{elsarticle}

%% Use the options 1p,twocolumn; 3p; 3p,twocolumn; 5p; or 5p,twocolumn
%% for a journal layout:
%% \documentclass[final,1p,times,authoryear]{elsarticle}
%% \documentclass[final,1p,times,twocolumn,authoryear]{elsarticle}
%% \documentclass[final,3p,times,authoryear]{elsarticle}
%% \documentclass[final,3p,times,twocolumn,authoryear]{elsarticle}
%% \documentclass[final,5p,times,authoryear]{elsarticle}
%% \documentclass[final,5p,times,twocolumn,authoryear]{elsarticle}

%% For including figures, graphicx.sty has been loaded in
%% elsarticle.cls. If you prefer to use the old commands
%% please give \usepackage{epsfig}

%% The amssymb package provides various useful mathematical symbols
\usepackage[utf8]{inputenc}
\usepackage{amsmath}
\usepackage{amscd}
\usepackage{amsthm}
\usepackage{ upgreek }
\usepackage{amssymb}
\usepackage{hyperref}
\usepackage{enumitem}
\usepackage{url}
\usepackage{times}
\newtheorem{theorem}{Theorem}
\newtheorem{lemma}{Lemma}[section]
\newtheorem{remark}{Remark}
\numberwithin{equation}{section}

%% The amsthm package provides extended theorem environments
%% \usepackage{amsthm}

%% The lineno packages adds line numbers. Start line numbering with
%% \begin{linenumbers}, end it with \end{linenumbers}. Or switch it on
%% for the whole article with \linenumbers.
%% \usepackage{lineno}

\journal{Journal of Number Theory}
\begin{document}
	\begin{frontmatter}
	
	%% Title, authors and addresses
	
	%% use the tnoteref command within \title for footnotes;
	%% use the tnotetext command for theassociated footnote;
	%% use the fnref command within \author or \affiliation for footnotes;
	%% use the fntext command for theassociated footnote;
	%% use the corref command within \author for corresponding author footnotes;
	%% use the cortext command for theassociated footnote;
	%% use the ead command for the email address,
	%% and the form \ead[url] for the home page:
	%% \title{Title\tnoteref{label1}}
	%% \tnotetext[label1]{}
	%% \author{Name\corref{cor1}\fnref{label2}}
	%% \ead{email address}
	%% \ead[url]{home page}
	%% \fntext[label2]{}
	%% \cortext[cor1]{}
	%% \affiliation{organization={},
		%%            addressline={}, 
		%%            city={},
		%%            postcode={}, 
		%%            state={},
		%%            country={}}
	%% \fntext[label3]{}
	
	
	\title{On the discrete mean square of certain hybrid sum involving $a_{\mathbb{K}}(n)$}
	
	%% use optional labels to link authors explicitly to addresses:
	
	\author[label1]{EKTA SONI}
	\author[label2]{M.S. DATT}
	\author[label3]{AYYADURAI SANKARANARAYANAN}
	\affiliation[label1]{organization={University of Hyderabad},%Department and Organization
		addressline={CR Rao Road, Gachibowli}, 
		city={Hyderabad},
		postcode={500046}, 
		state={Telangana},
		country={India}}
	\affiliation[label2]{organization={University of Hyderabad},%Department and Organization
		addressline={CR Rao Road, Gachibowli}, 
		city={Hyderabad},
		postcode={500046}, 
		state={Telangana},
		country={India}}
		\affiliation[label3]{organization={University of Hyderabad},%Department and Organization
			addressline={CR Rao Road, Gachibowli}, 
			city={Hyderabad},
			postcode={500046}, 
			state={Telangana},
			country={India}}
	\begin{abstract}
	%% Text of abstract
	Let $\mathbb{K}$ be a non-normal algebraic number field of cubic degree given by the polynomial $x^{3}+ax^{2}+bx+c$ of discriminant $D_{\mathbb{K}}<0$. For sufficiently large $x$, we establish an asymptotic formula for the hybrid sum
$$\sum\limits_{\substack{n= \sum_{i=1}^{8}n_{i}^{2}\leq x\\ (n_{1},n_{2},n_{3},n_{4},n_{5},n_{6},n_{7},n_{8})\in \mathbb{Z}^{8}  }} a_{\mathbb{K}}^{2}(n)$$
with a tight error term.
\end{abstract}

%%Graphical abstract
%\begin{graphicalabstract}
%\includegraphics{grabs}
%\end{graphicalabstract}

%%Research highlights
%\begin{highlights}
%\item Research highlight 1
%\item Research highlight 2
%\end{highlights}

\begin{keyword}
	Dedekind zeta function, Hecke eigenforms, Normalized Fourier coefficient, Symmetric square $L$-function, Dirichlet character.
	%% keywords here, in the form: keyword \sep keyword
	
	%% PACS codes here, in the form: \PACS code \sep code
	
	%% MSC codes here, in the form: \MSC code \sep code
	\MSC[2000] 11M \sep 11M06 \sep 11F11 \sep 11F30\sep 11R42.
	
\end{keyword}

\end{frontmatter}

%% \linenumbers
	\section{Introduction}
Let $\mathbb{K}$ be a number field of degree $m$ over $\mathbb{Q}$. Let $\mathcal{O}_{\mathbb{K}}$ be the ring of integers in $\mathbb{K}$. The Dedekind zeta function associated to $\mathbb{K}$ is defined as 
\begin{center}
	$\zeta_{\mathbb{K}}(s)=\prod\limits_{\substack{\mathcal{P}\subseteq \mathcal{O}_{\mathbb{K}}\\ \mathcal{P} \neq (0)}}\big(1-\frac{1}{N(\mathcal{P})}\big)^{-1}=\sum\limits_{\substack{\mathfrak{a}\subseteq \mathcal{O}_{\mathbb{K}}\\\mathfrak{a}\neq (0)}}\frac{1}{N^{s}(\mathfrak{a})}$,
\end{center}
for any $s\in \mathbb{C}$, $Re(s)>1$, where the product is taken over all non-zero prime ideals in $\mathcal{O}_{\mathbb{K}}$ and $N(\mathfrak{a})$ is the absolute norm of non-zero integral ideals.\\
The Dedekind zeta function can be extended meromorphically to the whole complex plane with a simple pole at $Re(s)=1$. The residue of the Dedekind zeta function at $s=1$ is given by
\begin{center}
	$\frac{2^{r_{1}}(2\pi)^{r_{2}}hR}{\omega \sqrt{|D_{\mathbb{K}}|}}$,
\end{center} 
where $r_{1}$ is the total number of real embeddings and $r_{2}$ is the total number of complex embeddings, $h$ is the class number, $R$ denotes the regulator, $\omega$ is the number of units in $\mathcal{O}_{\mathbb{K}}$ and $|D_{\mathbb{K}}|$ is the absolute value of the discriminant of $\mathbb{K}$.\\

%Let
%\begin{center} $\zeta_{\mathbb{K}}(s)=\Big(\frac{\sqrt{D_{\mathbb{K}}}}{2^{r_{2}}\pi^{\frac{n}{2}}}\Big)^{5}\Gamma(\frac{s}{2})^{r}\Gamma(s)^{r}\zeta_{\mathbb{K}}(s)$.
%\end{center}
%The Dirichlet series $\zeta_{\mathbb{K}}$ is analytic on the entire complex plane except for the simple pole at $s=0$ and $s=1$. And $\zeta_{\mathbb{K}}$ satisfies the functional equation 
%\begin{center}
%$\zeta_{\mathbb{K}}(s)=\zeta_{\mathbb{K}}(s-1)$.
%\end{center}
For $n\in \mathbb{N}$, let $a_{\mathbb{K}}(n)$ be the number of integral ideals of norm $n$ in $\mathcal{O}_{\mathbb{K}}$, then Dedekind zeta function $\zeta_{\mathbb{K}}(s)$ can be written as
\begin{center}
	$\zeta_{\mathbb{K}}(s)=\sum\limits_{n=1}^{\infty}\frac{a_{\mathbb{K}}(n)}{n^{s}},~~~~~Re(s)>1.$
\end{center}
As shown by Chandrasekharan and Good \cite{CKAG}, the coefficients $a_{\mathbb{K}}(n)$ are multiplicative and satisfy the following upper bound
%The coefficients $a_{\mathbb{K}}(n)$ defines a multiplicative arithmetical function $i.e,$ if $gcd(n,m)=1$, then $a_{\mathbb{K}}(mn)=a_{\mathbb{K}}(m)a_{\mathbb{K}}(n)$.\\
%By \cite{CKAG}, for any $\epsilon>0$,
\begin{center}
	$a_{\mathbb{K}}(n)\leq d(n)^{[\mathbb{K}:\mathbb{Q}]}\ll n^{\epsilon}$,
\end{center}
where $d(n)$ is the divisor function. As $a_{\mathbb{K}}(n)$ is multiplicative, we have
\begin{center}
	$\displaystyle \zeta_{\mathbb{K}}(s)=\prod\limits_{p}\Big(1+\frac{a_{\mathbb{K}}(p)}{p^{s}}+\frac{a_{\mathbb{K}}(p^{2})}{p^{2s}}+\dots+\frac{a_{\mathbb{K}}(p^{\ell})}{p^{\ell s}}+\dots\Big)$,
\end{center}
for $Re(s)>1$.\\
In $1949$, Landau\cite{Land}, in $1963$ Chandrasekharan and Narasimhan \cite{CKN} , and in $1983$ chandrasekharan and Good \cite{CKAG} considered first, second and $\ell^{th}$-moments $(\ell \geq 2)$ of $a_{\mathbb{K}}(n)$'s respectively and obtained asymptotic formulae.\\

For a cubic non-normal extension $\mathbb{K}$ over $\mathbb{Q}$, given by an irreducible polynomial $x^{3}+ax^{2}+bx+c\in \mathbb{Z}[x]$ such that the discriminant $D_{\mathbb{K}}<0$,
recently Naveen and Prashant \cite{P tiwari} studied the Problem for $\ell^{th}$-order moments of coefficients of the Dedekind zeta function over certain sequence of natural numbers and obtained asymptotic formulae, for a cubic non-normal extension over $\mathbb{Q}$.\\  

In this paper, we study the second order moments of coefficients of Dedekind zeta function over certain sequence of natural numbers. More precisely, we would like to investigate the asymptotic behaviour of $\sum\limits_{\substack{ n= \sum_{i=1}^{8}n_{i}^{2}\leq x\\ (n_{1},\dots,n_{8})\in \mathbb{Z}^{8}  }} a_{\mathbb{K}}^{2}(n)$ for sufficiently large $x$ and establish an asymptotic formula with a tight error term.\\

\noindent Indeed we prove the following:\\
\noindent\textbf{Main Theorem}:	\textit{ For $x\geq x_{0}$, where $x_{0}$ is sufficiently large, and $\epsilon>0$ any small constant, we have
	\begin{align*} 
		\sum\limits_{\substack{n= \sum_{i=1}^{8}n_{i}^{2}\leq x\\ (n_{1},n_{2},n_{3},n_{4},n_{5},n_{6},n_{7},n_{8})\in \mathbb{Z}^{8}  }} a_{\mathbb{K}}^{2}(n)&=C x^{4} P_{1}(\log{x})+O(x^{\frac{198}{53}+\epsilon}),
	\end{align*}
	where $P_{1}(\log{x})$ is a polynomial of degree $1$ in $\log{x}$.
}\\

The paper is planned as follows. In section $2$, we introduce some preliminaries and important lemmas that will play crucial role in the proof of our result. In section $3$, the proof of the theorem is presented. Throughout the paper, $\epsilon$ denotes sufficiently small positive constant, not necessarily the same at each occurrence and the implied constants may depend on $\epsilon,\,f$ and $|D_{\mathbb{K}}|$ and effective.  
\section{Preliminaries and Some Lemmas}
%Let $L(s,f)$ be the $L$-function associated with a normalized Hecke eigenforms $f$ of even integral weight $k$ for the full modular group $SL(2,\mathbb{Z})$. We denote by $\lambda_{f}(n)$ the normalized Fourier coefficient of the expansion \\
Let $k$ be a positive even integer and $H_{k}(SL(2,\mathbb{Z}))$ be the set of all normalized Hecke eigenforms of weight $k$ for the full modular group $SL(2,\mathbb{Z})$. For a cusp form $f\in H_{k}(SL(2,\mathbb{Z}))$, let $\lambda_{f}(n)$ be the normalized $n^{th}$ coefficients of Fourier expansion of $f$ at infinity, i.e, 
\begin{align*}
	f(z)&=\sum_{n=1}^{\infty} \lambda_{f}(n)n^{\frac{k-1}{2}}e^{2\pi i nz},~~~~\text{for all}~z\in{\mathbb{H}},
\end{align*}
where $\mathbb{H}$ be the Poincar\'e upper Half-plane. These Fourier coefficients $\lambda_{f}(n)$ are the Hecke eigenvalues of $f$ and $\lambda_{f}(n)\in \mathbb{R}$ $\forall~~n.$ The $\lambda_{f}(n)$ satisfies the following relation. For $m,n\in \mathbb{N}$
\begin{align*}
	\lambda_{f}(m)\lambda_{f}(n)&=\sum\limits_{d|gcd(m,n)}\lambda_{f}\Big(\frac{mn}{d^{2}}\Big).
\end{align*}
Thus $\lambda_{f}(n)$ is a multiplicative function.\\
In 1974, P. Deligne \cite{PD} proved that for any prime $p$, there exist complex numbers $\alpha_{f}(p),\,\beta_{f}(p)$ such that
\begin{align*}
	\lambda_{f}(p)&=\alpha_{f}(p)+\beta_{f}(p),\\
	|\alpha_{f}(p)|&=|\beta_{f}(p)|=\alpha_{f}(p)\beta_{f}(p)=1.
\end{align*}
Given any Dirichlet's character $\chi$ of Modulo $N$, we have the Dirichlet's $L$-function, namely
\begin{center}
	$L(s,\chi)=\sum\limits_{n=1}^{\infty}\frac{\chi(n)}{n^{s}}~~~\text{for}~Re(s)>1$,
\end{center}
and
\begin{align*}
	L(s,sym^{j}f)&=\prod_{p}\Big(\prod\limits_{i=0}^{j}\Big(1-\frac{\alpha^{j-i}_{f}(p)\beta^{i}_{f}(p)}{p^{s}}\Big)\Big)^{-1}\\
	&=\sum_{n=1}^{\infty}\frac{\lambda_{sym^{j}f}(n)}{n^{s}}~~~\text{for}~Re(s)>1.
\end{align*}
Note that $\lambda_{sym^{j}f}(n)$ is also real valued multiplicative function. We also have the Deligne's bound $\lambda_{f}(n)\leq d(n)\leq n^{\epsilon}$, $d(n)$ is the divisor function.\\
We note that the $L$-function associated to $\lambda_{f}(n)$ is defined as
\begin{align*}
	L(s,f)&=\sum_{n=1}^{\infty}\frac{\lambda_{f}(n)}{n^{s}}, 
\end{align*}
for $Re(s)>1,$ where $\lambda_{f}(n)$ is the eigenvalue of Hecke operators $T_{n}$ and the symmetric square $L$-function is defined as
\begin{align*}
	L(s,sym^{2}f)&=\sum\limits_{n=1}^{\infty}\frac{\lambda_{sym^{2}f}(n)}{n^{s}}\\
	&=\prod_{p} \Big(1-\frac{\alpha^{2}(p)}{p^{s}}\Big)^{-1}\Big(1-\frac{\beta^{2}(p)}{p^{s}}\Big)^{-1}\Big(1-\frac{1}{p^{s}}\Big)^{-1},
\end{align*}
for $Re(s)>1,$ where $\lambda_{sym^{2}f}(n)$ is multiplicative.\\
Let $N\in \mathbb{N}$ and $\chi$ be a Dirichlet character $(\mod N)$. Let $\Gamma_{0}(N)=\Big\{ \begin{pmatrix}
	a & b\\
	c & d
\end{pmatrix}
\in SL(2,\mathbb{Z})~\big|~\text{$c$ is a multiple of $N$}\Big\}$. The subgroup $\Gamma_{0}(N)$ is called the congruence subgroup of $SL(2,\mathbb{Z})$. For $\begin{pmatrix}
	a & b\\
	c & d
\end{pmatrix} \in \Gamma_{0}(N) $, and $z\in \mathbb{H}$ (the upper half plane), if $f \begin{pmatrix}
	az+b\\
	cz+d
\end{pmatrix}=\chi(d)(cz+d)^{k}f(z)$, then $f$ is known as the modular form of weight $k$ of level $N$ with Nebentypus $\chi$. \\
From \cite{P tiwari}, for a holomorphic cusp form  $f$ of weight $1$ of the congruence subgroup $\Gamma_{0}(N)$ with $N=|D_{\mathbb{K}}|$, we have 
\begin{center}
	$\zeta_{\mathbb{K}}(s)=\zeta(s)L(s,f),$
\end{center}
and hence,
\begin{center}
	$a_{\mathbb{K}}(n)=\sum\limits_{d|n}\lambda_{f}(d).$
\end{center}
This implies that $a_{\mathbb{K}}(p)=1+\lambda_{f}(p).$\\

\noindent Let
\begin{center}
	$	q(x_{1},x_{2},x_{3},x_{4},x_{5},x_{6},x_{7},x_{8})=\sum\limits_{i=1}^{8}x_{i}^{2}\in \mathbb{Z}[x_{1},x_{2},x_{3},x_{4},x_{5},x_{6},x_{7},x_{8}],$
\end{center} 
and
\begin{center}
	$r_{8}(n)=\#\{(a_{1},a_{2},a_{3},a_{4},a_{5},a_{6},a_{7},a_{8})\in \mathbb{Z}^{8}~|~ \sum_{i=1}^{8} a_{i}^{2}=n\}. $
\end{center}
From \cite{Jacobi}, we know that $r_{8}(n)=16(-1)^{n}\sum_{d|n}(-1)^{d}d^{3}$. First, we show that $(-1)^{n}\sum_{d|n}(-1)^{d}d^{3}$ is a multiplicative function.
\begin{lemma}
	Let $g(n)=(-1)^{n}\sum_{d|n}(-1)^{d}d^{3}$. Then $g(n)$ is a multiplicative arithmetic function.
\end{lemma}
\begin{proof}
	To prove the lemma, we show that for any two distinct primes $p,\,q$ and for $\ell,m \in \mathbb{N}$
	\begin{center}
		$g(p^{\ell}q^{m})=g(p^{\ell})g(q^{m})$.
	\end{center}
	\textbf{Case(i)} In this case we take $p=2$ and $q$ is odd. By the definition of $g(n)$,
	\begin{align*}
		g(2^{\ell})&=(-1)^{2^{\ell}}\{-1+2^{3}+2^{6}+\dots+2^{3\ell}\}\\
		&=\{-1+2^{3}+2^{6}+\dots+2^{3\ell}\},
	\end{align*}
	and 
	\begin{align*}
		g(q^{m})&=(-1)^{q^{m}}\{-1-q^{3}-\dots-q^{3m}\}\\
		&=\{1+q^{3}+\dots+q^{3m}\}.
	\end{align*}
	Therefore,
	\begin{align*}
		g(2^{\ell})g(q^{m})&=\{-1+2^{3}+2^{6}+\dots+2^{3\ell}\}\{1+q^{3}+\dots+q^{3m}\}\\
		&=-\{1+q^{3}+\dots+q^{3m}\}+\sum\limits_{\substack{1\leq i\leq \ell \\ 0\leq j \leq m}} 2^{3i} q^{3j}\\
		&=-\sum\limits_{j=0}^{m}q^{3j}+\sum\limits_{\substack{1\leq i\leq \ell \\ 0\leq j \leq m}} 2^{3i} q^{3j}.
	\end{align*}
	Now take,
	\begin{align*}
		g(2^{\ell}q^{m})&=(-1)^{2^{\ell}q^{m}}\big\{-1+\sum\limits_{i=1}^{\ell}2^{3i}-\sum\limits_{j=1}^{m}q^{3j}+\sum\limits_{\substack{1\leq i\leq \ell \\ 1\leq j \leq m}} 2^{3i} q^{3j}\big\}\\
		&=-\big\{1+\sum\limits_{j=1}^{m}q^{3j}\big\}+\big\{\sum\limits_{i=1}^{\ell}2^{3i}+ \sum\limits_{\substack{1\leq i\leq \ell \\ 1\leq j \leq m}} 2^{3i} q^{3j}\big\}\\
		&=-\sum\limits_{j=0}^{m}q^{3j}+\sum\limits_{\substack{1\leq i\leq \ell \\ 0\leq j \leq m}} 2^{3i} q^{3j}.
	\end{align*}
	Therefore, $g(2^{\ell}q^{m})=g(2^{\ell})g(q^{m})$.\vspace{3mm}\\
	\textbf{Case (ii)} Suppose both $p,\,q$ are odd primes.\\
	We have,
	\begin{align*}
		g(p^{\ell})&=(-1)^{p^{\ell}}\{-1-p^{3}-p^{6}\dots-p^{3\ell}\}\\
		&=1+p^{3}+p^{6}+\dots+p^{3\ell},
	\end{align*}
	and 
	\begin{align*}
		g(p^{m})&=1+q^{3}+q^{6}+\dots+q^{3m}.
	\end{align*}
	Therefore,
	\begin{align*}
		g(p^{\ell})g(p^{m})&=\{1+p^{3}+p^{6}+\dots+p^{3\ell}\}\{1+q^{3}+q^{6}+\dots+q^{3m}\}\\
		&=\sum\limits_{\substack{0\leq i\leq \ell \\ 0\leq j \leq m}} p^{3i}q^{3j}.
	\end{align*}
	Now, 
	\begin{align*}
		g(p^{\ell}q^{m})=(-1)^{p^{\ell}q^{m}}\sum\limits_{\substack{0\leq i\leq \ell \\ 0\leq j \leq m}}(-1)^{p^{3i}q^{3j}} p^{3i}q^{3j}=\sum\limits_{\substack{0\leq i\leq \ell \\ 0\leq j \leq m}} p^{3i}q^{3j}.
	\end{align*}
	Therefore,
	$g(p^{\ell}q^{m})=g(p^{\ell})g(q^{m})$.\\
	Hence the Lemma.
\end{proof}


We know that for instance by \cite{OM} , $a_{\mathbb{K}}(n)$ is multiplicative and by the above Lemma $g(n)$ is multiplicative. Hence $a_{\mathbb{K}}^{2}(n)g(n)$ is also multiplicative. To study the asymptotic behaviour of $\sum\limits_{\substack{ n= \sum_{i=1}^{8}n_{i}^{2}\leq x\\ (n_{1},n_{2},n_{3},n_{4},n_{5},n_{6},n_{7},n_{8})\in \mathbb{Z}^{8}  }}a_{\mathbb{K}}^{2}(n)$, we consider the Dirichlet's series associated to $a_{\mathbb{K}}^{2}(n)$ over the sequence of natural numbers given by $r_{8}(n)$. For any $x\in \mathbb{R^{+}},~x\geq 1$, we have
\begin{align*}
	\sum\limits_{\substack{n= \sum_{i=1}^{8}n_{i}^{2}\leq x\\ (n_{1},n_{2},n_{3},n_{4},n_{5},n_{6},n_{7},n_{8})\in \mathbb{Z}^{8}  }}a_{\mathbb{K}}^{2}(n) &=\sum_{n\leq x}a_{\mathbb{K}}^{2}(n)\Big(\sum\limits_{\substack{n= \sum_{i=1}^{8}n_{i}^{2}\\ (n_{1},n_{2},n_{3},n_{4},n_{5},n_{6},n_{7},n_{8})\in \mathbb{Z}^{8}}} 1\Big)\\
	&= \sum_{n\leq x} a_{\mathbb{K}}^{2}(n) r_{8}(n)\\
	&=16\sum_{n\leq x} a_{\mathbb{K}}^{2}(n) g(n).
\end{align*}

\begin{lemma} The $L$-function $F(s)$ associated with $a_{\mathbb{K}}^{2}(n)r_{8}(n)$ $i.e,$
	\begin{align*}
		F(s):=L(s, a_{\mathbb{K}}^{2}(n)r_{8}(n))&=\sum_{n=1}^{\infty}\frac{a_{\mathbb{K}}^{2}(n)r_{8}(n)}{n^{s}},
	\end{align*}
	can be expressed as
	\begin{align*}
		F(s)&=16\frac{B_{2}(s)}{A_{2}(s)}\zeta^{2}(s-3)L^{2}(s-3,f)L(s-3,sym^{2}f)\zeta^{2}(s)L^{2}(s,f)L(s,sym^{2}f)B(s).
	\end{align*}
	Let 
	\begin{center}
		$G(s)=16\zeta^{2}(s)L^{2}(s,f)L(s,sym^{2}f)B(s),$
	\end{center}
	where $B(s)$ is a harmless factor and $G(s)$ converges absolutely for $Re(s)>\frac{7}{2}$.
\end{lemma}
\begin{proof} 
	As $a_{\mathbb{K}}^{2}(n)$ and $(-1)^{n}\sum_{d|n}(-1)^{d}d^{3}$ are multiplicative, we have
	\begin{align*}
		F(s) &= L(s,a_{\mathbb{K}}^{2}(n)r_{8}(n))=16L(s,a_{\mathbb{K}}^{2}(n)g(n))\\
		&=16\sum_{n=1}^{\infty} \frac{a_{\mathbb{K}}^{2}(n)\big((-1)^{n}\sum_{d|n}(-1)^{d}d^{3}\big)}{n^{s}}\\
		&=16\prod_{p}\Big(1-\frac{a_{\mathbb{K}}^{2}(p)((-1)^{p+1}+p^{3})}{p^{s}}\Big)^{-1}.
	\end{align*}
	For $p\neq 2$, we let
	\begin{align*}
		A_{p}(s)&=1+\frac{a_{\mathbb{K}}^{2}(p)g(p)}{p^{s}}+\frac{a_{\mathbb{K}}^{2}(p^{2})g(p^{2})}{p^{2s}}+\cdots\\
		&=1+\frac{a_{\mathbb{K}}^{2}(p)(1+p^{3})}{p^{s}}+\frac{a_{\mathbb{K}}^{2}(p^{2})(1+p^{3}+p^{6})}{p^{2s}}+\cdots.
	\end{align*} 
	For $p=2$, we have
	\begin{align*}
		A_{2}(s)&=1+\frac{a_{\mathbb{K}}^{2}(2)(1+2^{3})}{2^{s}}+\frac{a_{\mathbb{K}}^{2}(2^{2})(1+2^{3}+2^{6})}{2^{2s}}+\cdots,
	\end{align*}
	and let
	\begin{align*}
		B_{2}(s)&=1+\frac{a_{\mathbb{K}}^{2}(2)(-1+2^{3})}{2^{s}}+\frac{a_{\mathbb{K}}^{2}(2^{2})(-1+2^{3}+2^{6})}{2^{2s}}+\cdots.
	\end{align*}
	Then,
	\begin{align*}
		F(s)&=16 B_{2}(s)\prod_{p\neq 2} A_{p}(s)\\
		&=16\frac{B_{2}(s)}{A_{2}(s)}\prod_{p} A_{p}(s) \\
		&=16\frac{B_{2}(s)}{A_{2}(s)}\prod_{p}\Big(1+\frac{a_{\mathbb{K}}^{2}(p)(1+p^{3})}{p^{s}}+\frac{a_{\mathbb{K}}^{2}(p^{2})(1+p^{3}+p^{6})}{p^{2s}}+\dots\Big)\\
		&=16\frac{B_{2}(s)}{A_{2}(s)}\prod_{p}\Big(1+\frac{(1+\lambda_{f}(p))^{2}(1+p^{3})}{p^{s}}+\frac{(1+\lambda_{f}(p^{2}))^{2}(1+p^{3}+p^{6})}{p^{2s}}+\dots\Big)\\
		&=16\frac{B_{2}(s)}{A_{2}(s)}\prod_{p}\Big(1+\frac{(2+2\lambda_{f}(p)+\lambda_{sym^{2}f}(p))(1+p^{3})}{p^{s}}+\dots\Big)\\
		&=16\frac{B_{2}(s)}{A_{2}(s)}\prod_{p}\Big(1+\frac{2+2\lambda_{f}(p)+\lambda_{sym^{2}f}(p)+2p^{3}+2p^{3}\lambda_{f}(p)+p^{3}\lambda_{sym^{2}f}(p)}{p^{s}}+\dots\Big)\\
		&=16\frac{B_{2}(s)}{A_{2}(s)}\zeta^{2}(s-3)L^{2}(s-3,f)L(s-3,sym^{2}f)\zeta^{2}(s)L^{2}(s,f)L(s,sym^{2}f)B(s),
	\end{align*}
	where $B(s)$ is a harmless factor which is absolutely convergent for $Re(s)>\frac{7}{2}$. The Dirichlet series $F(s)$ has a pole of order $2$ at $s=4$.\\
	
	\newpage
	\noindent Let
	\begin{center}
		$G(s)=16\zeta^{2}(s)L^{2}(s,f)L(s,sym^{2}f)B(s)$.
	\end{center}
	Now, we can write\\
	\begin{center}
		$F(s)=\frac{B_{2}(s)}{A_{2}(s)}\zeta^{2}(s-3)L^{2}(s-3,f)L(s-3,sym^{2}f)G(s)$.
	\end{center}
	We would like to show that $0<\Big|\frac{B_{2}(s)}{A_{2}(s)}\Big|\leq 1$ for $Re(s) =\frac{7}{2}$.\\
	As $\mathbb{K}/\mathbb{Q}$ is a field extension of degree $3$, the ideal generated by a prime number $p$ in $\mathbb{Z}$ can be expressed as a product of prime ideals in $\mathcal{O}_{\mathbb{K}}$ as
	\begin{center}
		$p\mathcal{O}_{\mathbb{K}}=\mathcal{P}_{1}^{e_{1}}\,\mathcal{P}_{2}^{e_{2}}\dots\mathcal{P}_{r}^{e_{r}}$,
	\end{center}
	with $e_{1}f_{1}+e_{2}f_{2}+\dots+e_{r}f_{r}=3$, where $e_{i}$ is the ramification index and $f_{i}$ is the degree of $\mathcal{O}_{\mathbb{K}}/\mathcal{P}_{i}$ over the field $\mathbb{Z}_{p}$ (see, Page.no. $52$, \cite{milne}). This implies that $r\leq 3$ and $p\mathcal{O}_{\mathbb{K}}=\mathcal{P}_{1}^{e_{1}}\,\mathcal{P}_{2}^{e_{2}}\,\mathcal{P}_{3}^{e_{3}}$.\\
	In particular for the ideal $(2)$ generated by $2$
	\begin{center}
		$(2)\mathcal{O}_{\mathbb{K}}=\mathcal{P}_{1}^{e_{1}}\,\mathcal{P}_{2}^{e_{2}}\,\mathcal{P}_{3}^{e_{3}}$.
	\end{center}
	We have the following three cases:\\
	\textbf{Case (i):} If $(2)\mathcal{O}_{\mathbb{K}}$ is a prime ideal in $\mathcal{O}_{\mathbb{K}}$, then $a_{\mathbb{K}}(2^{n})\leq1$.\\
	
	\noindent\textbf{Case (ii):} If $(2)\mathcal{O}_{\mathbb{K}}=\mathcal{P}_{1}\mathcal{P}_{2}\mathcal{P}_{3}$, $\mathcal{P}_{1},\,\mathcal{P}_{2},\,\mathcal{P}_{3}$ are distinct then $a_{\mathbb{K}}(2^{n})\leq 1$.\\
	
	\noindent\textbf{Case (iii):} If $(2)\mathcal{O}_{\mathbb{K}}=\mathcal{P}_{1}^{e_{1}}\mathcal{P}_{2}^{e_{2}}$, then $e_{1}+e_{2}=3$. Assume that $e_{1}=1$ and $e_{2}=2$ then $N((2)\mathcal{O}_{\mathbb{K}})=2^{3}$. Hence we have, if    $I$ is an ideal of $\mathcal{O}_{\mathbb{K}}$ such that $2^{n}=N(I)$, then  $I=\mathcal{P}_{1}^{e_{1}}\mathcal{P}_{2}^{e_{2}}$ such that $e_{1}+2e_{2}=3$. Therefore the number of integral ideals in $\mathcal{O}_{\mathbb{K}}$ whose norm is $2^{n}$, is less than or equal to $\Big[\frac{n}{2}\Big]+1\leq n$, for $n\geq 2.$\\
	
	\noindent For $s=\frac{7}{2}$, we show that $A_{2}(\frac{7}{2})$ converges absolutely. 
	\begin{align*}
		A_{2}(\frac{7}{2})&=1+\frac{a_{\mathbb{K}}^{2}(2)(2^{3}+1)}{2^{\frac{7}{2}}}+\frac{a_{\mathbb{K}}^{2}(2^{2})(2^{6}+2^{3}+1)}{2^{2.\frac{7}{2}}}+\cdots\\
		&=1+\frac{a_{\mathbb{K}}^{2}(2)(2^{6}-1)}{(2^{3}-1)2^{\frac{7}{2}}}+\frac{a_{\mathbb{K}}^{2}(2^{2})(2^{9}-1)}{(2^{3}-1)2^{2.\frac{7}{2}}}+\cdots\\
		&\leq 1+\frac{1}{7}\Big\{\frac{a_{\mathbb{K}}^{2}(2)2^{6}}{2^{\frac{7}{2}}}+\frac{a_{\mathbb{K}}^{2}(2^{2})2^{9}}{2^{2.\frac{7}{2}}}+\frac{a_{\mathbb{K}}^{2}(2^{3})2^{12}}{2^{3.\frac{7}{2}}}+\cdots\Big\}\\
		&=1+\frac{8}{7}\Big\{\frac{a_{\mathbb{K}}^{2}(2)}{2^{\frac{1}{2}}}+\frac{a_{\mathbb{K}}^{2}(2^{2})}{2}+\frac{a_{\mathbb{K}}^{2}(2^{3})}{2^{\frac{3}{2}}}+\cdots\Big\}\\
		&\leq 1+\frac{8}{7}\Big\{\frac{a_{\mathbb{K}}^{2}(2)}{\sqrt{2}}+\frac{a_{\mathbb{K}}^{2}(2^{2})}{(\sqrt{2})^{2}}+\frac{a_{\mathbb{K}}^{2}(2^{3})}{(\sqrt{2})^{3}}+\cdots\Big\}.
	\end{align*}
	We know that, $a_{\mathbb{K}}(2^{n})\leq n$, therefore $(\sqrt{2})^{n}=2^{\frac{n}{2}}=e^{\frac{n\log{2}}{2}}>\frac{(\frac{n\log{2}}{2})^{4}}{4!}$, which implies that $\frac{1}{(\sqrt{2})^{n}}<\frac{2.4!}{(n\log{2})^{4}}$.\\
	Hence,
	\begin{align*}
		\sum_{n=0}^{\infty} \frac{a_{\mathbb{K}}^{2}(2^{n})}{(\sqrt{2})^{n}}&< 1+\frac{8}{7}\sum_{n=1}^{\infty}\frac{a_{\mathbb{K}}^{2}(2^{n}).2.4!}{(n\log{2})^{4}}\\
		&\leq 1+\frac{8}{7}\sum_{n=1}^{\infty} \frac{(n)^{2}.2.4!}{n^{4}(\log{2})^{4}}\\
		&=1+\frac{2^{4}.4!}{7(\log 2)^{4}}\sum_{n=1}^{\infty} \frac{1}{n^{2}}\\
		&\leq\frac{2^{4}.4!}{7(\log 2)^{4}}\Big(1+\sum_{n=1}^{\infty}\frac{1}{n^{2}}\Big).
	\end{align*}
	Hence $1+\sum_{n=0}^{\infty} \frac{a_{\mathbb{K}}^{2}(2^{n})}{(\sqrt{2})^{n}}$ is a convergent series, and hence $A_{2}(\frac{7}{2})$ converges.\\
	As, $\frac{a_{\mathbb{K}}^{2}(2^{n})(2^{3n}+2^{3(n-1)}+\dots+2^{3}-1)}{2^{\frac{7}{2}n}}\leq\frac{a_{\mathbb{K}}^{2}(2^{n})(2^{3n}+2^{3(n-1)}+\dots+2^{3}+1)}{2^{\frac{7}{2}n}}$ for every $n\in \mathbb{N}$, and $\frac{a_{\mathbb{K}}^{2}(2^{n})(2^{3n}+2^{3(n-1)}+\dots+2^{3}-1)}{2^{\frac{7}{2}n}}>0$, and hence $B_{2}(\frac{7}{2})$ also converges to real number $>0$. Therefore by the Theorem $11.4$ \cite{apostol}, $0<|B_{2}(s)|<|A_{2}(s)|$ for $Re(s)>\frac{7}{2}$ and hence $0<\big|\frac{B_{2}(s)}{A_{2}(s)}\big|\leq 1$ for $Re(s)>\frac{7}{2}$ .
	This implies that, $\frac{B_{2}(s)}{A_{2}(s)}G(s)\neq 0$, and is absolutely convergent for $Re(s)>\frac{7}{2}$.\\
	The Dirichlet's series $F(s)$ has a pole of order $2$ at $s=4$ and converges absolutely for $Re(s)>4$. It does not vanish on the line $Re(s)=4.$
\end{proof}

In the course of the proof of the theorem, we shall make use of certain well-known estimates for the Riemann zeta function and associated $L$-functions, which we state below. 
\begin{lemma}\cite{Kram}
	\textit{	For $U\geq U_{0},$ where $U_{0}$ is sufficiently large, there exists a point $T^*\in (U,\,2U)$ such that
		$$\max\limits_{\sigma \geq \frac{1}{2}}|\zeta(\sigma\pm iT^*)|~\leq~exp(C(\log\log{U})^{2})\ll U^{\epsilon}.$$}
\end{lemma}
\begin{lemma} \cite{Ylin}
	\textit{	For any $\epsilon>0,$ we have
		\begin{align*}
			\zeta(\sigma+it)&\ll_{\epsilon} (1+|t|)^{max\{\frac{13}{42}(1-\sigma),0\}+\epsilon},
		\end{align*}
		uniformly for $\frac{1}{2}\leq \sigma \leq 1+\epsilon$ and $|t|\geq 1.$}
\end{lemma}
\begin{lemma}\cite{Aivik}
	\textit{For any $\epsilon>0$, we have 
		\begin{align*}
			L(\sigma+it,f)&\ll (|t|+10)^{max\{\frac{2}{3}(1-\sigma),0\}+\epsilon},
		\end{align*}
		uniformly for $\frac{1}{2}\leq \sigma\leq 1+\epsilon$ and $|t|\geq 10$.}
\end{lemma}
\begin{lemma}\cite{Ylin}
	\textit{For any $\epsilon>0$, we have
		\begin{align*}
			L(\sigma+it,sym^{2}f)&\ll (|t|+10)^{max\{\frac{6}{5}(1-\sigma),0\}+\epsilon},
		\end{align*}
		uniformly for $\frac{1}{2}\leq \sigma\leq 1+\epsilon$ and $|t|\geq 10$.}
\end{lemma}
\begin{lemma}\cite{AS}
	\textit{For $T\geq 100$, we have 
		\begin{align*}
			\int_{10}^{T}|L(\frac{1}{2}+\epsilon+it,sym^{2}f)|^{2}\,dt&\ll T^{\frac{3}{2}+\epsilon}.
		\end{align*}
	}
\end{lemma}

\begin{lemma}\cite{Lsankar}
	\textit{
		For $T\geq 100$, we have 
		\begin{align*}
			\int_{10}^{T}|L(\frac{1}{2}+\epsilon+it,f)|^{2}\,dt&\ll T^{1+\epsilon}.
		\end{align*}
	}
\end{lemma}
\section{Proof of the main Theorem}
To start with, we make the choice of our $T=\pm T^{\star}\in (U,2U)$ to satisfy Lemma $2.3$. By applying Perron's formula \cite{Agranville} to $F(s)$ with the choice $\eta= 4+\epsilon$ and $1\leq T \leq x$, we get
\begin{align*}
	\sum\limits_{\substack{n\leq x \\
			n=n_{1}^{2}+n_{2}^{2}+\dots+n_{8}^{2}\\(n_1,n_2,\dots,n_{8})\in \mathbb{Z}^{8}}} a_{\mathbb{K}}^{2}(n)&= \frac{1}{2\pi i} \int_{4+\epsilon-iT}^{4+\epsilon+iT} F(s) \frac{x^{s}}{s}\,ds+O\Big(\frac{x^{4+\epsilon}}{T}\Big).
\end{align*}
By moving the line of integration to $Re(s)=\frac{7}{2}+\epsilon$. We consider the rectangle $R$ with the vertices $\frac{7}{2}+\epsilon\pm iT$ and $4+\epsilon\pm iT$. In the rectangle $R$, $F(s)\frac{x^{s}}{s}$ has a pole of multiplicity $2$ at $s=4$ coming from the factor $\zeta^{2}(s-3)$. Thus, by using Cauchy's residue theorem on $R$, we get
\begin{align*}
	\frac{1}{2\pi i}\int_{4+\epsilon-iT}^{4+\epsilon+iT}F(s)\frac{x^{s}}{s}\,ds&=M(x)+\frac{1}{2\pi i}\Big\{-\int_{4+\epsilon+iT}^{\frac{7}{2}+\epsilon+iT}-\int_{\frac{7}{2}+\epsilon+iT}^{\frac{7}{2}+\epsilon-iT}\\
	&~~~~~~~~~~~~~~~~~~~~~~~-\int_{\frac{7}{2}+\epsilon-iT}^{4+\epsilon-iT}\Big\}F(s) \frac{x^{s}}{s}\,ds\\
	&=M(x)+\frac{1}{2\pi i}(J_{1}+J_{2}+J_{3})~~~~~~~(\text{say}),
\end{align*}
here $M(x)$ is the main term i.e, the residue of $F(s)\frac{x^{s}}{s}$ at $s=4$, and is given by $Cx^{4}P_{1}(\log{x})$ where $P_{1}(\log{x})$ is a polynomial of degree $1$ in $\log{x}$.\\

Let $C_{H}$ be the horizontal contribution of $F(s)\frac{x^{s}}{s}$ along the lines $\frac{7}{2}+\epsilon-iT$ to $4+\epsilon-iT$ and $\frac{7}{2}+\epsilon+iT$ to $4+\epsilon+iT$.
\begin{align*}
	|C_{H}|& \ll \Bigg|\int_{\frac{7}{2}+\epsilon+iT}^{4+\epsilon+iT}\,F(s)\frac{x^{s}}{s}\,ds\Bigg|\\
	&\ll \int_{\frac{7}{2}+\epsilon+iT}^{4+\epsilon+iT} |\zeta(s-3)|^{2}|L(s-3,f)|^{2}|L(s-3,sym^{2}f)|\Big|\frac{A_{2}(s)}{B_{2}(s)}\Big||G(s)|\frac{|x^{s}|}{|s|}\,ds.
	%	&\ll \frac{x^{3}}{T}\int_{\frac{1}{2}+\epsilon}^{1+\epsilon}|\zeta(\sigma+iT)|^{2}|L(\sigma+iT,f)|^{2}|L(\sigma+iT,sym^{2}f)|\,x^{\sigma}\,d\sigma.
\end{align*}
As $\Big|\frac{A_{2}(s)}{B_{2}(s)}\Big||G(s)|$ is bounded for $Re(s)>\frac{7}{2}$, we have
\begin{align*}
	|C_{H}|&\ll\frac{x^{3}}{T}\int_{\frac{1}{2}+\epsilon}^{1+\epsilon}|\zeta(\sigma+iT)|^{2}|L(\sigma+iT,f)|^{2}|L(\sigma+iT,sym^{2}f)|\,x^{\sigma}\,d\sigma.
\end{align*}
By using Lemmas $2.3$, $2.5$ and $2.6$ we get
\begin{align*}
	|C_{H}|& \ll \frac{x^{3}}{T}\,T^{2\epsilon}\int_{\frac{1}{2}+\epsilon}^{1+\epsilon} T^{\frac{4}{3}(1-\sigma)+\epsilon}\,T^{\frac{6}{5}(1-\sigma)+\epsilon}\,x^{\sigma}\,d\sigma\\
	&\ll x^{3}T^{\frac{23}{15}+4\epsilon}\int_{\frac{1}{2}+\epsilon}^{1+\epsilon} \Big(\frac{x}{T^{\frac{38}{15}}}\Big)^{\sigma}\, d\sigma.
\end{align*}
Since $\Big(\frac{x}{T^{\frac{38}{15}}}\Big)^{\sigma}$ is a monotonic function of $\sigma$ on $[\frac{1}{2}+\epsilon,1+\epsilon]$, it follows that the maximum is attained at the endpoints of the interval $[\frac{1}{2}+\epsilon,1+\epsilon]$. Thus,
\begin{align*}
	|C_{H}|&\ll x^{3}T^{\frac{23}{15}+4\epsilon} \Big(\Big(\frac{x}{T^{\frac{38}{15}}}\Big)^{1+\epsilon}+\Big(\frac{x}{T^{\frac{38}{15}}}\Big)^{\frac{1}{2}+\epsilon}\Big)\\
	& \ll \frac{x^{4+\epsilon}}{T}+x^{\frac{7}{2}+\epsilon}T^{\frac{4}{15}+\epsilon}.
\end{align*}
Now we estimate the contribution along the left vertical line $(J_{2})$ in absolute value
\begin{align*}
	|C_{V}|&=\Big|\int_{\frac{7}{2}+\epsilon-iT}^{\frac{7}{2}+\epsilon+iT} F(s)\frac{x^{s}}{s}\,ds\Big|\\
	&=\int_{\frac{7}{2}+\epsilon-iT}^{\frac{7}{2}+\epsilon+iT} |\zeta(s-3)|^{2}|L(s-3,f)|^{2}|L(s-3,sym^{2}f)|\Big|\frac{A_{2}(s)}{B_{2}(s)}\Big||G(s)|\frac{|x^{s}|}{|s|}\,ds\\
	&\ll x^{\frac{7}{2}+\epsilon}+x^{\frac{7}{2}+\epsilon}\int_{10}^{T} |\zeta(\frac{1}{2}+\epsilon+it)|^{2}|L(\frac{1}{2}+\epsilon+it,f)|^{2}|L(\frac{1}{2}+\epsilon+it,sym^{2}f)|\frac{1}{t}\,dt\\
	&\ll x^{\frac{7}{2}+\epsilon}+x^{\frac{7}{2}+\epsilon} \max_{10\leq t\leq T}\{|\zeta(\frac{1}{2}+\epsilon+it)|^{2}|L(\frac{1}{2}+\epsilon+it,f)|\}\\
	&~~~~~~~~~~\times \int_{10}^{T} |L(\frac{1}{2}+\epsilon+it,f)|\,|L(\frac{1}{2}+\epsilon+it,sym^{2}f)|\frac{1}{t}\,dt.
\end{align*}
By using Lemmas $2.4$, $2.5$, $2.7$ and $2.8$ we get
\begin{align*}
	|C_{V}| &\ll x^{\frac{7}{2}+\epsilon}+ x^{\frac{7}{2}+\epsilon}\,(T^{\frac{13}{84}})^{2}\,T^{\frac{1}{3}+\epsilon}\Big(\int_{10}^{T}\frac{|L(\frac{1}{2}+\epsilon+it,f)|^{2}}{t}\,dt\Big)^{\frac{1}{2}}\Big(\int_{10}^{T}\frac{|L(\frac{1}{2}+\epsilon+it,sym^{2}f)|^{2}}{t}\,dt\Big)^{\frac{1}{2}}\\
	& \ll x^{\frac{7}{2}+\epsilon}+ x^{\frac{7}{2}+\epsilon}\,T^{\frac{9}{14}+\epsilon} T^{\frac{1}{4}+\epsilon} \\
	&\ll x^{\frac{7}{2}+\epsilon}\,T^{\frac{25}{28}+\epsilon}.
\end{align*}
Therefore, the total error is (note that $\frac{4}{15}<\frac{25}{28}$)
\begin{align*}
	|C_{H}|+|C_{V}|+O\Big(\frac{x^{4+\epsilon}}{T}\Big)&\ll\frac{x^{4+\epsilon}}{T}+x^{\frac{7}{2}+\epsilon}T^{\frac{25}{28}+\epsilon}+O\Big(\frac{x^{4+\epsilon}}{T}\Big).
\end{align*}
To get an optimal estimate, we choose $T=x^{\frac{14}{53}}$ so that we obtain
\begin{align*}
	|C_{H}|+|C_{V}|+O\Big(\frac{x^{4+\epsilon}}{T}\Big)&= O(x^{\frac{198}{53}+\epsilon}).
\end{align*}
Therefore, 
\begin{align*}
	\sum\limits_{\substack{n=\sum_{i=1}^{8}n_{i}^{2} \leq x\\(n_1,n_2,n_{3},n_{4},n_{5},n_{6},n_{7},n_{8})\in \mathbb{Z}^{8}}} a_{\mathbb{K}}^{2}(n)&= Cx^{4}P_{1}(\log{x})+O(x^{\frac{198}{53}+\epsilon}).
\end{align*}
This proves the main Theorem.
\hfill $\Box$
\begin{thebibliography}{99}
	\bibitem{apostol} Apostol, Tom M. Introduction to analytic number theory. Book \emph{Springer Science and Business Media}, 2013.
	\bibitem{CKAG} Chandrasekharan, K., and A. Good, On the number of integral ideals in Galois extensions. \emph{onatshefte für Mathematik}, 95.2 (1983): 99-109.
	\bibitem{CKN} Chandrasekharan, K., and Raghavan Narasimhan, The approximate functional equation for a class of zeta-functions, \emph{Mathematische Annalen}, 152.1 (1963): 30-64.
	\bibitem{PD} Deligne, Pierre, La conjecture de Weil. I, \emph{ Publications Mathématiques de l'Institut des Hautes Études Scientifiques}, 43.1 (1974) 273-307.
	\bibitem{OM} Fomenko, O.M., Mean values associated with the Dedekind zeta function. \emph{Zap. Nauchn. Sem. S.-Peterburg. Otdel. Mat. Inst. Steklov.} (POMI), 350:187-198, 203, (2007).
	\bibitem{Agranville} A. Granville and K. Soundararajan, Multiplication number theory: The pretentious approach. Book \emph{manuscript in preparation}, 12, 2014.
	\bibitem{hardy} Hardy, Godfrey Harold, and Edward Maitland Wright, An introduction to the theory of numbers. \emph{Oxford university press}, 1979.
	\bibitem{Aivik} Ivic, Aleksander, The Riemann Zeta-function: The theory of the Riemann zeta-functions with applications, \emph{John Wiley and Sons, New York}, 1985.
	\bibitem{Jacobi} Jacobi, Carl Gustav Jacob. Fundamenta nova theoriae functionum ellipticarum.
	\emph{Cambridge University Press}, 2012.
	\bibitem{Land} Landau, Edmund: Einfuhrung in die elementare und analytische Theorie der algebraischen Zahlen und der Ideale. \emph{Chelsea Publishing Co., New York}, (1949).
	\bibitem{Lsankar} Lao, HuiXue, and Ayyadurai Sankaranarayanan, On the mean-square of the error term related to $ \sum_{n\leq x} \lambda^{2}(n^{j})$, \emph{Sci. China Math.}, 54.5 (2011) 855-864.
	\bibitem{Ylin} Y. Lin, R. Nunes and Z. Qi, Strong subconvexity for self-dual $GL(3)$ $L$-functions, \emph{Int. Math. Res. Not.}, 3(2023) 11453-11470. 0
	\bibitem{milne} Milne, James S., Algebraic number theory. 2020.
	\bibitem{P tiwari} Naveen K. Godara and Prashant Tiwari, The higher power moments of the coefficients of the Dedekind zeta function over a polynomial in six variables, \emph{Integers}, 25(2025).
	\bibitem{Kram} K. Ramachandra and A. Sankaranarayanan, Notes on the Riemann zeta-function. \emph{J. Indian Math. Soc.}, {\bf 57} (1991), 67-77.
	\bibitem{AS} A.Sankaranarayanan, Fundamental properties of symmetric square $ L $-functions. I, \emph{lllinois Journal of Mathematics}, 46.1 (2002) 23-43.
\end{thebibliography}
\end{document}
